% Diagramma originale; licenza CC BY-SA 4.0 del materiale didattico.
\begin{tikzpicture}[unipd/diagram]
  \draw[unipd/edge,rounded corners=2pt] (-4.6,0.5) rectangle (4.6,2.7);
  \node at (0,2.4) {CPU};
  \node[unipd/architecture component] (control) at (-3,1.35)
    {Unità di\\controllo};
  \node[unipd/architecture component] (alu) at (0,1.35) {ALU};
  \node[unipd/architecture component] (registers) at (3,1.35)
    {Registri\\PC, istruzione\\dati};
  \draw[unipd/edge,<->] (control) -- (alu);
  \draw[unipd/edge,<->] (alu) -- (registers);
  \draw[unipd/edge,line width=1.2pt] (-4.6,-0.3) -- (4.6,-0.3);
  \node[font=\normalfont,anchor=north] at (0,-0.4)
    {Bus: dati, indirizzi, controllo};
  \draw[unipd/edge,<->] (0,0.5) -- (0,-0.3);
  \node[unipd/architecture component] (memory) at (-3,-1.5)
    {Memoria centrale\\istruzioni e dati};
  \node[unipd/architecture component] (io) at (3,-1.5)
    {Interfacce I/O\\periferiche};
  \draw[unipd/edge,<->] (memory.north) -- (-3,-0.3);
  \draw[unipd/edge,<->] (io.north) -- (3,-0.3);
\end{tikzpicture}
